\section*{第 \ref{ch:g7:organs-at-work} 章——工作中的器官}

\begin{solution}{exo:g7:organs-at-work:1}
从\omterm{prop:g4:journey-of-food:absorb}{血液}中摄取的\omterm{def:g7:what-is-respiration:respiration}{氧气}和\omterm{prop:g7:organs-at-work:bill}{葡萄糖}增多；释放的
\omterm{def:g7:what-is-respiration:respiration}{二氧化碳}和热量增多。
\end{solution}

\begin{solution}{exo:g7:organs-at-work:2}
\omterm{prop:g4:journey-of-food:absorb}{血液}中现成的糖类养料——\omterm{def:g5:first-look-at-cells:cell}{细胞}燃烧时的
日常燃料。\omterm{def:g3:bones-and-muscles:muscle}{肌肉}从流经它的\omterm{prop:g4:journey-of-food:absorb}{血液}中摄取
这种养料（养料由上游的\omterm{def:g4:journey-of-food:digestion}{消化}过程
供应并储存在\omterm{prop:g4:journey-of-food:absorb}{血液}中）。
\end{solution}

\begin{solution}{exo:g7:organs-at-work:3}
将进入\omterm{def:g3:bones-and-muscles:muscle}{肌肉}的\omterm{prop:g4:journey-of-food:absorb}{血液}与离开\omterm{def:g3:bones-and-muscles:muscle}{肌肉}的\omterm{prop:g4:journey-of-food:absorb}{血液}进行比较。
静息时，离开的\omterm{prop:g4:journey-of-food:absorb}{血液}只释放了少量\omterm{def:g7:what-is-respiration:respiration}{氧气}和\omterm{prop:g7:organs-at-work:bill}{葡萄糖}；
剧烈运动时，离开的\omterm{prop:g4:journey-of-food:absorb}{血液}几乎释放了全部可输送的\omterm{def:g7:what-is-respiration:respiration}{氧气}，
\omterm{def:g7:what-is-respiration:respiration}{二氧化碳}则大大增多。
两者之差即\omterm{def:g3:bones-and-muscles:muscle}{肌肉}的“账单”。
\end{solution}

\begin{solution}{exo:g7:organs-at-work:4}
\omterm{def:g7:what-is-respiration:respiration}{呼吸}加深加快；\omterm{def:g4:heart-and-blood:heart}{心脏}跳动更快、更有力；
\omterm{prop:g4:journey-of-food:absorb}{血液}重新分配：工作\omterm{def:g3:bones-and-muscles:muscle}{肌肉}和散热的\omterm{def:g1:the-five-senses:sense}{皮肤}\omterm{def:g4:heart-and-blood:vessels}{血管}扩张，
静息时的\omterm{def:g4:journey-of-food:tract}{消化道}
\omterm{def:g4:heart-and-blood:vessels}{血管}则收缩。
\end{solution}

\begin{solution}{exo:g7:organs-at-work:5}
燃料燃烧总会释放热量，工作中的\omterm{def:g3:bones-and-muscles:muscle}{肌肉}
是人体的“熔炉”。打寒战就是\omterm{def:g3:bones-and-muscles:muscle}{肌肉}被驱动，
只为产生热量——按需供热。
\end{solution}

\begin{solution}{exo:g7:organs-at-work:6}
运动结束后，\omterm{prop:g4:heart-and-blood:pulse}{脉搏}和\omterm{def:g7:what-is-respiration:respiration}{呼吸}逐渐\omterm{def:g7:organs-at-work:recovery}{恢复}平稳。
身体正在补充\omterm{def:g7:what-is-respiration:respiration}{氧气}储备、清除运动产生的废物，
并散发积累的热量。
\end{solution}

\begin{solution}{exo:g7:organs-at-work:7}
静息时约 $72$；平台期接近 $160$；
图中\omterm{def:g7:organs-at-work:recovery}{恢复}时间约为五至六分钟。
训练会降低静息值，
降低相同运动量下的平台值，并缩短\omterm{def:g7:organs-at-work:recovery}{恢复}时间。
\end{solution}

\begin{solution}{exo:g7:organs-at-work:8}
两者都是\omterm{prop:g4:journey-of-food:absorb}{血液}重新分配造成的效应。午饭后，\omterm{def:g4:journey-of-food:tract}{消化道}占用了大量\omterm{prop:g4:journey-of-food:absorb}{血液}，
身体其他部位便显得迟缓；
此时冲刺会迫使两处争夺\omterm{prop:g4:journey-of-food:absorb}{血液}供应，
——岔气，以及“饭后不宜剧跑”的古老忠告，
都与此有关。
\end{solution}

\begin{solution}{exo:g7:organs-at-work:9}
\omterm{def:g1:the-five-senses:sense}{皮肤}\omterm{def:g4:heart-and-blood:vessels}{血管}扩张，是把热量从“熔炉”输送到体表的“散热器”：
向外散热，解决工作\omterm{def:g3:bones-and-muscles:muscle}{肌肉}过热的问题。
面色发红说明散热系统
正在运转，而非故障。
\end{solution}

\begin{solution}{exo:g7:organs-at-work:10}
\omterm{def:g4:heart-and-blood:heart}{心脏}每次搏动输送更多\omterm{prop:g4:journey-of-food:absorb}{血液}，这是日常可见的标志，
静息\omterm{prop:g4:heart-and-blood:pulse}{脉搏}更慢。\omterm{def:g3:bones-and-muscles:muscle}{肌肉}增粗，储存更多燃料，
同样的坡道走起来更轻松。供血网络的
细小\omterm{def:g4:heart-and-blood:vessels}{血管}增多，\omterm{def:g7:what-is-respiration:respiration}{呼吸}加深——运动时仍能说话，
\omterm{def:g7:organs-at-work:recovery}{恢复}也更快。
\end{solution}

\begin{solution}{exo:g7:organs-at-work:11}
静息 $68$、\omterm{def:g7:organs-at-work:recovery}{恢复}仅需三分钟的机器，正是受过训练的那一台。
静息\omterm{prop:g4:heart-and-blood:pulse}{脉搏}说明\omterm{def:g4:heart-and-blood:heart}{心脏}每次搏出更多\omterm{prop:g4:journey-of-food:absorb}{血液}；
\omterm{def:g7:organs-at-work:recovery}{恢复}时间则说明供血线路能迅速偿还运动欠下的“债务”。
二者正是
\cref{prop:g7:organs-at-work:training}所述重建的两端。
\end{solution}

\begin{solution}{exo:g7:organs-at-work:12}
冲刺时，腿部\omterm{def:g3:bones-and-muscles:muscle}{肌肉}\omterm{def:g5:first-look-at-cells:cell}{细胞}借助\omterm{def:g7:what-is-respiration:respiration}{氧气}燃烧\omterm{prop:g7:organs-at-work:bill}{葡萄糖}以提供收缩动力——
这正是\omterm{def:g7:what-is-respiration:respiration}{呼吸作用}的目的——而它们的“账单”
（燃料与\omterm{def:g7:what-is-respiration:respiration}{氧气}流入增多，\omterm{def:g7:what-is-respiration:respiration}{二氧化碳}与热量流出增多）必须立即结清。
\omterm{prop:g4:journey-of-food:absorb}{血液}充当结算服务的角色：
在肺和肠道装载养分与\omterm{def:g7:what-is-respiration:respiration}{氧气}，
由加速跳动的\omterm{def:g4:heart-and-blood:heart}{心脏}加快泵送，
再重新分配到腿部，输送燃料与\omterm{def:g7:what-is-respiration:respiration}{氧气}，运走\omterm{def:g7:what-is-respiration:respiration}{二氧化碳}——
一次冲刺，三章内容，同一笔账。
\end{solution}

\begin{solution}{pb:g7:organs-at-work:1}
\textbf{1.} \omterm{def:g3:bones-and-muscles:muscle}{肌肉}从流经的\omterm{prop:g4:journey-of-food:absorb}{血液}中摄取的\omterm{def:g7:what-is-respiration:respiration}{氧气}和\omterm{prop:g7:organs-at-work:bill}{葡萄糖}增多；
\omterm{def:g7:what-is-respiration:respiration}{二氧化碳}和热量的产生量随之上升——
这就是剧烈运动的“账单”。

\textbf{2.} 因为二者是同一供血链条上的两个环节：
\omterm{def:g7:what-is-respiration:respiration}{呼吸}在肺部为\omterm{prop:g4:journey-of-food:absorb}{血液}补充\omterm{def:g7:what-is-respiration:respiration}{氧气}，
\omterm{def:g4:heart-and-blood:heart}{心脏}再把补充后的\omterm{prop:g4:journey-of-food:absorb}{血液}输送到\omterm{def:g3:bones-and-muscles:muscle}{肌肉}。
一端的需求提高，两者的节奏也随之加快。

\textbf{3.} 工作中的\omterm{def:g3:bones-and-muscles:muscle}{肌肉}是“熔炉”，其产生的热量必须散出。
\omterm{prop:g4:journey-of-food:absorb}{血液}将热量带到\omterm{def:g1:the-five-senses:sense}{皮肤}，
\omterm{def:g1:the-five-senses:sense}{皮肤}\omterm{def:g4:heart-and-blood:vessels}{血管}扩张，面色发红并向外散热——
内部是熔炉，表面是散热器。

\textbf{4.} 因为仍有欠账：
需要补充\omterm{def:g7:what-is-respiration:respiration}{氧气}储备、清除废物、散发热量。
只有账目结清，\omterm{prop:g4:heart-and-blood:pulse}{脉搏}才会下降——这种下降就是\omterm{def:g7:organs-at-work:recovery}{恢复}。

\textbf{5.} 静息 $84 \to 72$：\omterm{def:g4:heart-and-blood:heart}{心脏}经过重建，
每次搏出血量增多。运动 $168 \to 150$：同样的台阶
现在负担更小——\omterm{def:g3:bones-and-muscles:muscle}{肌肉}更强壮，
细小\omterm{def:g4:heart-and-blood:vessels}{血管}供应更充足。\omterm{def:g7:organs-at-work:recovery}{恢复} $9 \to 5$ 分钟：整条链路
偿还欠账的速度加快。本学期的体育锻炼
完成了这次重建。

\textbf{6.} 要么 B 学期初就已受过训练——
其数值接近训练后的下限，变化很小——或者 B
跳过了本学期的体育锻炼，只是维持着一台稳定的
未经训练的机器。关键事实：B 在课外的
活动（若在别处有运动，则证实前一种解读；
若无，则证实后一种——B 第一周的良好数据
已倾向于此）。

\textbf{7.} 改道规则：进餐后，\omterm{def:g4:journey-of-food:tract}{消化道}
会抢占\omterm{prop:g4:journey-of-food:absorb}{血液}，因此静息\omterm{prop:g4:heart-and-blood:pulse}{脉搏}读数
偏高，运动耐受力偏低。测量方案：午餐前测量，
或餐后两小时以上测量——
且始终选择同一时段。

\textbf{8.} 公平的比较：只有以相同方式采集的测量值
才能相互比较，无论是周与周之间，
还是学生与学生之间。

\textbf{9.} “十二周提高了 A 的上限：同样的
运动现在只需 $150$ 次搏动，而非 $168$ 次，
静息\omterm{prop:g4:heart-and-blood:pulse}{脉搏}低了八次。而且 A 的机器
结清欠账只需五分钟，九月时则需九分钟——
\omterm{def:g7:organs-at-work:recovery}{恢复}是所有收获中最明显的。”

\textbf{10.} 通常恰恰相反，不必担心：受过训练的
\omterm{def:g4:heart-and-blood:heart}{心脏}每次搏动输出的\omterm{prop:g4:journey-of-food:absorb}{血液}更多，因此静息时
频率较低——骑行者平静时的 $50$ 次，
承担的是未经训练者 $75$ 次的工作。（如有担忧，
应交由医生判断，但静息心率低而体能好是常见情况。）

\textbf{11.} \omterm{def:g7:organs-at-work:recovery}{恢复}时间衡量的是整条链路——
\omterm{def:g4:heart-and-blood:heart}{心脏}、肺、\omterm{def:g4:heart-and-blood:vessels}{血管}、\omterm{def:g3:bones-and-muscles:muscle}{肌肉}——偿还真实的欠账，
并随训练稳定且明显地改善；
运动\omterm{prop:g4:heart-and-blood:pulse}{脉搏}则更多取决于当天的确切运动量和
状态。作为单一跟踪指标，\omterm{def:g7:organs-at-work:recovery}{恢复}时间
是更可靠的概括。

\textbf{12.} 其用途本身即是：运动，温和而
有规律地反复进行，正是它重建了执行运动的机器。
\end{solution}
